\chapter{Wymagania i narzędzia}
\label{ch:wymagania-i-narzedzia}

% \begin{itemize}
% \item wymagania funkcjonalne i niefunkcjonalne
% \item opis narzędzi, metod eksperymentalnych, metod modelowania itp.
% \item metodyka pracy nad projektowaniem i implementacją -- dla prac, w których ma to zastosowanie
% \end{itemize}

\section{Wymagania funkcjonalne}

\section{Wymagania niefunkcjonalne}

\section{Narzędzia}
\subsection{Narzędzia głowne}
\subsubsection{ANTLR} % metody eksperymentalne, rozwijanie gramatyk, generacja api antlera do .go
\subsubsection{go} % biblioteki?
\subsubsection{C} % w kontekscie pisania bibliotek do bbuss
\subsubsection{GCC} % asm -> bin, preprocesor
\subsubsection{Linux} % jako docelowa platforma programów generowanych przez kompilator

\subsection{Metody eksperymentalne}
\subsubsection{nasm}
\subsubsection{ld}

\subsection{Narzędzia do rozwóju oprogramowania}
\subsubsection{Środowisko do edycji tekstu} 
% nvim: go debugger, lsp
% terminal
% tmux
\subsubsection{git}
% kontrola wersji
